% Options for packages loaded elsewhere
\PassOptionsToPackage{unicode}{hyperref}
\PassOptionsToPackage{hyphens}{url}
\PassOptionsToPackage{dvipsnames,svgnames,x11names}{xcolor}
\documentclass[
  11pt,
  letterpaper,
]{article}
\usepackage{xcolor}
\usepackage[margin=1in]{geometry}
\usepackage{amsmath,amssymb}
\setcounter{secnumdepth}{-\maxdimen} % remove section numbering
\usepackage{iftex}
\usepackage[OT1]{fontenc}
\renewcommand{\rmdefault}{cmr}
\renewcommand{\sfdefault}{cmss}
\renewcommand{\ttdefault}{cmtt}

\ifPDFTeX\else
  % xetex/luatex font selection
\fi
% Use upquote if available, for straight quotes in verbatim environments
\IfFileExists{upquote.sty}{\usepackage{upquote}}{}
\IfFileExists{microtype.sty}{% use microtype if available
  \usepackage[]{microtype}
  \UseMicrotypeSet[protrusion]{basicmath} % disable protrusion for tt fonts
}{}
\makeatletter
\@ifundefined{KOMAClassName}{% if non-KOMA class
  \IfFileExists{parskip.sty}{%
    \usepackage{parskip}
  }{% else
    \setlength{\parindent}{0pt}
    \setlength{\parskip}{6pt plus 2pt minus 1pt}}
}{% if KOMA class
  \KOMAoptions{parskip=half}}
\makeatother
\setlength{\emergencystretch}{3em} % prevent overfull lines
\providecommand{\tightlist}{%
  \setlength{\itemsep}{0pt}\setlength{\parskip}{0pt}}
\usepackage{microtype}
\usepackage{graphicx}
\setlength{\emergencystretch}{3em}
\allowdisplaybreaks[1]
\usepackage{bookmark}
\IfFileExists{xurl.sty}{\usepackage{xurl}}{} % add URL line breaks if available
\urlstyle{same}
% fallback for those not using the hyperref driver hyperxmp:
\makeatletter
\@ifundefined{xmpquote}{\newcommand{\xmpquote}[1]{#1}}{}
\makeatother
\hypersetup{
  pdftitle={Research notes and remaining questions},
  colorlinks=true,
  linkcolor={Maroon},
  filecolor={Maroon},
  citecolor={Blue},
  urlcolor={blue},
  pdfcreator={LaTeX via pandoc}}

\title{Research notes and remaining questions}
\usepackage{etoolbox}
\makeatletter
\providecommand{\subtitle}[1]{% add subtitle to \maketitle
  \apptocmd{\@title}{\par {\large #1 \par}}{}{}
}
\makeatother
\subtitle{A guide to the current results and exploratory targets}
\author{}
\date{21 October 2026}

\begin{document}
\maketitle
\begin{abstract}
This companion distinguishes proved results, conditional applications
and questions recorded during the Alexander research project. The
complete pseudo-visibility, species and quitting arguments have their
own current papers. We give a self-contained finite-observation
obstruction and explain its conditional application to a companion
population construction. The remaining cards state the additional
semantic or quantitative choices required before an exploratory
suggestion becomes a theorem target. A recorded question is not a
certification of present-day openness or novelty.
\end{abstract}

\section{1. Where the full arguments are now
presented}\label{where-the-full-arguments-are-now-presented}

The project's results concern several different notions of learning,
representation and comparison. Keeping them in separate full papers
makes their assumptions easier to assess.

\textbf{Pseudo-visibility.} The full paper
\href{../12-pseudovisibility-private-information/pseudovisibility-private-information.pdf}{\emph{Information
and learning in pseudo-visibility environments}} separates deterministic
policy representation, private-query value, shared random seeds and
learning dynamics. It gives exact one-step values and a restricted
sample-average tabular convergence theorem. The included finite-horizon
experiments are evidence about specified learners, rather than a
universal wrapper-improvement theorem.

\textbf{Species.} The full paper
\href{../06-founding-cohort-species/founding-cohort-species.pdf}{\emph{Specieslike
clusters with several founders}} combines founding cohorts and fixed
founder bounds, proves coverage, compares infinite generator-equivalence
classes, and gives the Borel complexity and binary non-learning results
for IAP. The equivalence statement concerns infinite clusters. A theorem
about deciding IAP does not by itself rule out every notion of
convergence for species-valued outputs.

\textbf{Quitting.} The full paper
\href{../01-quitters-and-question-26/quitters.pdf}{\emph{Quitting
thresholds and reward normalization in universal intelligence}} combines
the ultrafilter boundary with pathwise and signed expected-value models.
For the latter, its counterexample has quantifiers
\(\forall r\,\exists U\,\exists A,B\). The universal reference machine
is chosen as part of the witness.

These are current arguments with their assumptions stated directly. This
companion does not provide alternative versions of their proofs.

\section{2. A finite-observation
obstruction}\label{a-finite-observation-obstruction}

Consider an infinite binary stream \(s\). A learner sees successive
finite prefixes and outputs a binary guess. Eventual correctness means
that its guesses stabilize to the right answer on every input; the
learner need not be computable.

\subsection{Proposition. Eventual periodicity has no universally correct
learner}\label{proposition.-eventual-periodicity-has-no-universally-correct-learner}

No finite-prefix learner is eventually correct on every stream about
whether that stream is eventually periodic.

\textbf{Proof.} Let EP be the eventually periodic streams in Cantor
space. It is countable: a finite initial segment and finite period
specify a member. It is dense: every finite prefix can be continued
periodically. Each singleton is nowhere dense, so EP is meager.

If the guesses \(g_n\), based on length-\(n\) prefixes, converged
correctly on every input, then \[
                     \mathrm{EP}
                        =\bigcup_N\bigcap_{n\geq N}\{g_n=1\}
\] would be \(F_\sigma\), since each guess set is clopen. The same
expression with zero would make its complement \(F_\sigma\), hence EP
would also be \(G_\delta\). But a dense \(G_\delta\) subset of Cantor
space is comeager. It cannot also be meager, by the Baire category
theorem. This is a contradiction. \hfill\(\square\)

The argument is topological; it does not use a limitation of computer
running time. Any task continuously encoding this yes/no property
inherits the obstruction.

The public companion
\href{https://github.com/avg-netizen/biological-unavoidability-results}{\emph{Biological
unavoidability results}} gives a population family \(P_s\) for which
realizing every infinite binary sequence is equivalent to eventual
periodicity of \(s\). Under that construction and its finite-prefix
observation interface, the proposition gives non-learning of its
universality property. The construction itself is a separately cited
input, not proved here.

Finite certification, one-sided detection and eventual two-sided
correctness are different requirements. The property that a 1 occurs is
detected at its first occurrence. An algorithm can start with no and
switch to yes then. A prefix containing a 1 certifies the positive
answer; an all-zero prefix cannot certify the negative one. Promises can
change a learning problem as well: a stream promised to repeat a
specified-length period after a known cutoff is determined by a
sufficiently long prefix.

\section{3. Pattern arithmetic: specify the representation
first}\label{pattern-arithmetic-specify-the-representation-first}

Alexander's \emph{Arithmetical algorithms for elementary patterns}
(Archive for Mathematical Logic, 2015,
\href{https://doi.org/10.1007/s00153-014-0404-9}{DOI:
10.1007/s00153-014-0404-9}), Section 6, proposes further algorithms for
epsilon, Gamma and Veblen functions, and for second-order patterns.

A useful algorithmic target must fix the pattern class, coding of
inputs, notational range, operation and correctness interpretation. For
example, an epsilon-function algorithm would take a pattern denoting
\(\alpha\), return one denoting \(\varepsilon_\alpha\), and establish
correctness under the chosen pattern/ordinal correspondence. A
two-argument Veblen algorithm needs an equally explicit domain and
convention for \(\varphi_\alpha(\beta)\).

These are recorded specifications, not claims that arbitrary pattern
arithmetic is elementary, that small tests prove correctness, or that
the literature has left each task open. A new selection requires
checking the precise existing algorithm and current source version.

\section{4. Variadic Apply: the function argument is part of the
problem}\label{variadic-apply-the-function-argument-is-part-of-the-problem}

Alexander's \emph{The first-order syntax of variadic functions} (Notre
Dame Journal of Formal Logic 54(1), 2013,
\href{https://arxiv.org/abs/1105.4135}{arXiv:1105.4135}), Section 7,
motivates an Apply operation: \[
                \operatorname{Apply}(G,n_1,\ldots,n_k)
                           =G(n_1,\ldots,n_k).
\] Giving a sort for finite sequences does not by itself make the
variable function \(G\) available.

An exact target must choose a function sort, a code domain with an
evaluation relation, or an indexed operator family. It must say whether
the functions are all intended functions or a specified represented
collection. Its completeness claim must distinguish ordinary first-order
models from any intended full function semantics. Conservativity also
needs an explicit original language and class of reducts.

Until those choices are made, the card is a semantic design question
rather than a proved completeness theorem.

\section{5. Questions retained with their
limits}\label{questions-retained-with-their-limits}

\subsection{Internodon membership and
persistence}\label{internodon-membership-and-persistence}

The
\href{../08-internodon-order-robustness/internodon-order-robustness.pdf}{internodon
paper} proves a clade-change bound and states a separate merge-tree
height-stability theorem. Neither directly controls the co-membership of
fixed organisms.

A possible next question is whether an additional persistence margin,
together with an organism-label correspondence and the SD/test-organism
rules, guarantees membership stability. The separating branch and margin
must first be defined for that labelled construction. Ordinary diagram
stability alone does not supply this implication. The saved all-pairs
percentages do not establish it either.

\subsection{Intelligence sign for a fixed symmetric
machine}\label{intelligence-sign-for-a-fixed-symmetric-machine}

Alexander--Hutter's \emph{Reward-punishment symmetric universal
intelligence} (2021,
\href{https://doi.org/10.1007/978-3-030-93758-4_1}{DOI:
10.1007/978-3-030-93758-4\_1}), Section 5, raises sign computability.
Its duality gives zero to reward-ignoring agents.

A recorded reduction strategy would run a reward-ignoring agent until a
program halts, then switch to one of finitely many reward-sensitive
rules. To decide halting from the strong sign, one needs a uniform
non-vanishing statement: for every switch time and every reference
machine under consideration, at least one switched rule has nonzero
value. That statement is not supplied here. Adding the reward dual of a
rule only negates its value and does not repair vanishing. The weak-sign
variant has its own promise. Neither the quitting paper nor finite
special-prior constructions settle this fixed-machine question.

\subsection{Introspection outside the unstratified
lemma}\label{introspection-outside-the-unstratified-lemma}

Entry
\href{../07-minimal-introspection-bundle/minimal-introspection-bundle.pdf}{07}
analyzes the unstratified introspection bundle. The stratified
comparison concerns Lemma 65 and Corollary 71 of \emph{Self-referential
theories} (2020). A proposed extension must check where compactness,
assigned validity and closure operate, and whether closure only moves
upward in the ordinal index. Reusing the unstratified proof's wording
does not establish that transfer. This card records the dependency
rather than claiming a new minimality theorem.

\subsection{Credit assignment and learning
rates}\label{credit-assignment-and-learning-rates}

The pseudo-visibility paper proves equal greedy limits only for its
finite tabular model with sample averages, unique maximizers and
explicit sampling/noise assumptions. Extending that result to all
algorithms that update erased parameters from erased rewards is a
broader question. A bound on the gap, or on constant versus
context-linear sample complexity, needs specified rewards, learner,
features, accuracy and confidence. The existing experiment does not pay
those requirements.

\subsection{Representation and music
complexity}\label{representation-and-music-complexity}

Alexander's playful \emph{Legg--Hutter universal intelligence implies
classical music is better than pop music for intellectual training} (The
Reasoner, 2019) invites a representation-level comparison. A score, a
recording and a performance conditional on a score are different coding
objects. A meaningful comparison should hold that level fixed and state
the compressor, transcription and conditioning information.

This is an exploratory empirical suggestion, not a proved genre ranking
or a conclusion about Kolmogorov complexity from file sizes. No
experiment for it is included.

\section{6. How to use the cards}\label{how-to-use-the-cards}

Each card separates a specified result from the missing semantic,
quantitative or source-status input. The purpose is to make a next
question assessable. A saved conjecture does not establish present-day
openness, and a finite control does not replace an infinite proof.

The full papers linked in Section 1 provide their own sources,
assumptions and reproducibility guides. The remaining cards record
possibilities without claiming a literature survey, research priority or
new resolution.

\end{document}
